\section{Model}

\subsection{Environment}

The economy consists of a continuum of agents indexed by $i \in [0,1]$. Time is static. There is an underlying fundamental $\theta \in \mathbb{R}$ representing exchange rate misalignment, where higher values of $\theta$ correspond to a greater overvaluation of the official exchange rate.

The fundamental $\theta$ is drawn from a normal prior distribution:
\begin{equation}
\theta \sim \mathcal{N}(\mu, \sigma_\theta^2).
\end{equation}

Agents are risk averse and derive utility from consumption. They face uncertainty about the true level of exchange rate misalignment and must form beliefs based on imperfect information.

\subsection{Motives for Foreign Currency Demand}

Agents demand foreign currency for two distinct reasons:

\begin{enumerate}
\item \textbf{Hedging motive:} Agents acquire foreign currency to insure against the risk of a devaluation.
\item \textbf{Speculative motive:} Agents acquire foreign currency to exit ahead of a potential regime collapse.
\end{enumerate}

These two motives generate different components of aggregate foreign currency demand, which play distinct roles in equilibrium.

\subsection{Assets}

Agents can access foreign currency through two instruments:

\paragraph{Cash foreign currency ($c$):}
\begin{itemize}
\item Involves higher transaction costs
\item Trades in decentralized and opaque markets
\item Provides only noisy, privately observed information
\end{itemize}

\paragraph{Stablecoins ($s$):}
\begin{itemize}
\item Lower transaction costs
\item Scalable and accessible through digital platforms
\item Generates a publicly observable market price
\end{itemize}

Let $q_c$ and $q_s$ denote holdings of cash and stablecoins, respectively. Total foreign currency holdings are given by:
\begin{equation}
q = q_c + q_s.
\end{equation}

\subsection{Transaction Costs}

The cost of acquiring foreign currency depends on the instrument used. We assume convex transaction costs:
\begin{align}
C_c(q_c) &= a_c + b_c \sqrt{q_c} + \kappa_c, \\
C_s(q_s) &= a_s + b_s \sqrt{q_s} + \kappa_s,
\end{align}

with $C_s(q_s) < C_c(q_c)$ for comparable quantities, reflecting the efficiency advantage of stablecoins.

\subsection{Discussion}

The two-asset structure is essential. It allows us to separate the transaction cost channel from the information channel. A one-asset model with improved information would confound these mechanisms. In contrast, the presence of both cash and stablecoins enables us to study how financial innovation affects both market efficiency and the information environment.
